%% ============================================================================
%% Section: Background and Problem
%% ============================================================================
\section{Background and Problem Statement}\label{sec:bg}

In this section, we introduce the preliminary notions needed as basic
ingredients in our framework: causal DAGs and their property-graph encoding,
the catalog of structural queries, and the data-engineering problem we
address.

\subsection{Causal DAGs and their property-graph encoding}\label{sec:bg:dag}

\begin{definition}[Causal DAG]
A causal DAG is a directed acyclic graph $G=(V,E)$ where $V$ is a set of
random variables and $E \subseteq V \times V$ is a set of edges such that
$X \rightarrow Y \in E$ asserts that $X$ is a direct cause of $Y$ relative to
the other variables in $V$~\cite{pearl2009causality}. A function
$\lambda: V \rightarrow \{\textit{observed}, \textit{latent}\}$ assigns to
each variable whether it is measured.
\end{definition}

Latent variables are retained in the model because they change which
conditional independencies hold and which effects are identifiable, even
though no data is ever assigned to them. We write $\pa(X)$, $\ch(X)$, $\an(X)$,
and $\de(X)$ for the parents, children, ancestors, and descendants of $X$;
ancestors and descendants are the transitive closures of parents and children
and exclude $X$ itself. For instance, in Fig.~\ref{fig:example}
$\pa(\textit{onset\_covid19}) = \{\textit{exposure\_sarscov2},
\textit{access\_tests}, \textit{age}, \textit{susceptibility}\}$ and
$\an(\textit{reported\_cases})$ contains every drawn variable except
\textit{reported\_cases} itself.

\noindent\textbf{Property-graph encoding.} Causal4j stores the DAG as a native
property graph. Each variable is a node
$(:\textsf{Variable}\;\{\textsf{id},\ \textsf{name},\ \textsf{observed},\ \textsf{latent}\})$
and each direct cause is a relationship
$(:\textsf{Variable})\text{-}[:\textsf{CAUSES}\;\{\textsf{active},\ \textsf{weight}\}]\!\rightarrow\!(:\textsf{Variable})$.
The \textsf{id} property is the canonical key; \textsf{active} lets a curator
toggle an edge without deleting it, and only edges with
$\textsf{active}=\text{true}$ participate in the model; \textsf{weight}
carries an optional structural coefficient used by the quantitative layer
(Section~\ref{sec:quant}) and ignored by the structural queries. This
encoding keeps the causal model in the database, co-located and jointly
queryable with any other graph data, rather than in a transient in-memory
object owned by an external tool.

\subsection{The causal query catalog}\label{sec:bg:catalog}

Analysts interrogate the DAG through a small and well-established catalog of
structural queries, which we group by the reasoning they support.

\noindent\textbf{Reachability queries} return $\pa$, $\ch$, $\an$, and $\de$
of a variable, and test whether a directed path $X \rightarrow \cdots
\rightarrow Y$ exists. They underpin every other query.

\noindent\textbf{Path-role queries} classify the variables that lie between a
treatment $T$ and an outcome $Y$. A \textit{mediator} lies on a directed path
$T \rightarrow \cdots \rightarrow Y$, i.e., it belongs to
$\de(T) \cap \an(Y)$; a \textit{confounder} is a common cause,
i.e., a variable in $\an(T) \cap \an(Y)$; a \textit{collider} between $X$ and
$Y$ is a variable $Z \notin \an(X) \cup \an(Y)$ with one parent in
$\de(X) \cup \{X\}$ and one parent in $\de(Y) \cup \{Y\}$, i.e., a variable
at which a path from $X$ and a path from $Y$ meet head to head. In
Fig.~\ref{fig:example}, \textit{covid19\_burden} is a confounder of
$T$ and $Y$, \textit{exposure\_sarscov2} and \textit{onset\_covid19} are
mediators, and \textit{onset\_covid19} is a collider between \textit{age} and
\textit{susceptibility}.

\noindent\textbf{Independence queries} test d-separation.

\begin{definition}[d-separation]
Let $X, Y \in V$ and $Z \subseteq V \setminus \{X,Y\}$. A path between $X$ and
$Y$ is blocked by $Z$ if it contains a non-collider that belongs to $Z$, or a
collider such that neither the collider nor any of its descendants belongs to
$Z$. $X$ and $Y$ are d-separated given $Z$, written $X \dsep Y \mid Z$, if
every path between them is blocked by $Z$~\cite{pearl2009causality}.
\end{definition}

D-separation is the graphical criterion for conditional independence and is
not plain reachability: adding a variable to $Z$ can block one path and, when
that variable is a collider or one of its descendants, open another. In
Fig.~\ref{fig:example}, $\textit{age} \dsep \textit{susceptibility} \mid
\emptyset$ holds, while conditioning on \textit{reported\_cases} opens the
path through the collider \textit{onset\_covid19}.

\noindent\textbf{Identification queries} decide whether the causal effect of
$T$ on $Y$ can be estimated from observational data by covariate adjustment
and, if so, return an adjustment set. We adopt the generalized adjustment
criterion~\cite{shpitser2010validity,perkovic2018complete}, which is complete
for DAGs and is the criterion implemented by reference
tools~\cite{textor2016dagitty}.

\begin{definition}[Adjustment criterion]
Let $T, Y \in V$ with $Y \in \de(T)$. The causal nodes are
$\cn(T,Y) = (\de(T) \cap \an(Y)) \cup \{Y\}$, i.e., the non-treatment
variables on directed paths from $T$ to $Y$, and the forbidden set is
$\forb(T,Y) = \de(\cn(T,Y)) \cup \cn(T,Y) \cup \{T\}$. The proper backdoor
graph $G^{\mathit{pbd}}_{T,Y}$ is $G$ without the edges $T \rightarrow C$ for
$C \in \cn(T,Y)$. A set $Z \subseteq V$ satisfies the adjustment criterion
relative to $(T,Y)$ if $Z \cap \forb(T,Y) = \emptyset$ and
$T \dsep Y \mid Z$ holds in $G^{\mathit{pbd}}_{T,Y}$.
\end{definition}

The criterion has a constructive form that we exploit in
Section~\ref{sec:index}: the effect is identifiable by adjustment if and only
if the canonical set $\adj(T,Y) = \an(\{T,Y\}) \setminus (\forb(T,Y) \cup
\{T,Y\})$ satisfies the criterion~\cite{vanderzander2019separators}. In
Fig.~\ref{fig:example}, $\{\textit{covid19\_burden}\}$ satisfies the
criterion for $(\textit{interventions}, \textit{reported\_cases})$, while
\textit{mobility} is forbidden since it lies on a directed path from $T$ to
$Y$. Listing all inclusion-minimal sets that satisfy the criterion is the most
expensive query of the catalog, since each candidate is validated by a
d-separation test.

\subsection{Problem statement}\label{sec:bg:problem}

We target the following data-engineering problem: supporting the causal query
catalog as a first-class workload inside a graph DBMS. Concretely, a solution
has to meet four requirements.
\begin{itemize}
  \item \textbf{R1 (Expressiveness).} A declarative interface in which each
  catalog query is a single statement, close to how analysts phrase causal
  questions, that compiles to physical operators over the graph store.
  \item \textbf{R2 (Efficiency).} Structural queries answered without
  re-traversing the stored graph on every call, so that the reachability and
  path-role operators are sublinear in graph size and identification issues a
  bounded number of conditioning-aware traversals.
  \item \textbf{R3 (Correctness).} Results faithful to the graphical semantics
  above, in particular a conditioning-aware d-separation test and an
  adjustment criterion that is complete for DAGs.
  \item \textbf{R4 (Maintainability).} The accelerating structure kept
  consistent under edge and variable updates without full recomputation, since
  curated causal graphs change continually.
\end{itemize}
None of these is served by exporting the DAG to an external library: R2 and
R4 in particular are database concerns (indexing, query processing, and
incremental maintenance), which is why we address them inside the DBMS.
